\documentclass[11pt]{article}
\usepackage[margin=1.05in]{geometry}
\usepackage{amsmath,amssymb,amsthm,booktabs,longtable}
\usepackage[hidelinks]{hyperref}
\usepackage[T1]{fontenc}
\usepackage{lmodern}
\hypersetup{pdftitle={Minimum central circles: an effective characterization of the global asymptotic constant},
 pdfauthor={Maurizio Falconi}}
\newtheorem{theorem}{Theorem}
\newtheorem{lemma}[theorem]{Lemma}
\newtheorem{proposition}[theorem]{Proposition}
\newtheorem{corollary}[theorem]{Corollary}
\theoremstyle{definition}\newtheorem{definition}[theorem]{Definition}
\theoremstyle{remark}\newtheorem{remark}[theorem]{Remark}
\newcommand{\Rstar}{R^\ast}
\newcommand{\ch}{R_{\mathrm{chain}}}
\newcommand{\full}{R_{\mathrm{full}}}
\newcommand{\CC}{C_\ast}
\DeclareMathOperator{\atan}{arctan}
\DeclareMathOperator{\asin}{arcsin}
\title{Minimum central circles:\\an effective characterization of the global asymptotic constant}
\author{Maurizio Falconi}
\date{September 11, 2026}
\begin{document}
\maketitle
\begin{abstract}
Let $\Rstar(n)$ be the least radius of a central circle to which nonoverlapping
circles of radii $1,\ldots,n$ are externally tangent. We prove that
$\Rstar(n)=\CC n^2+o(n^2)$ and characterize $\CC$ by finite linear programs
with an explicit error tending to zero. The reduction preserves arbitrary
orders and all pairwise constraints: the limiting problem places marked
points on a line at pairwise separation at least the geometric mean of their
marks. Concatenation with a bounded boundary cost proves existence, and
balanced finite-word programs supply matching effective upper and lower
bounds. Their certified gap is $(1/k+1/r)/\pi$, before directed arithmetic
error, for $k$ mark types and words of length $r$.
A quantitative reflected-block recovery theorem supplies genuine permutations
and full ring geometry, with a countable extension and a strict four-block
improvement. The explicit interval is
$C_{\rm term}+\eta_{\rm width}\le\CC\le U_4$; neither endpoint is asserted sharp.
In particular, the earlier proposed coefficient $1/8$ is false.
An elementary expression for $\CC$, efficient high-precision evaluation and
global floating-circle structure remain open.
\end{abstract}

\section{Problem and relation to prior finite work}

All outer circles are tangent to the central circle; outer circles need
not be tangent to each other. For a cyclic order $\sigma$, the adjacent
chain radius $\ch(\sigma)$ ignores nonadjacent constraints, whereas
$\full(\sigma)$ is the least radius satisfying every pair constraint in
that order. The global quantity is $\Rstar(n)=\min_\sigma\full(\sigma)$.
Always $\ch(\sigma)\le\full(\sigma)$. A feasible construction supplies
an upper bound, not an optimum certificate.

The preceding study \cite{v1} established the anti-Monge/Supnick connection
for adjacent-chain ordering and developed the all-pairs finite optimization
and numerical certification framework. It reported certified global optima
for $3\le n\le14$, with an absolute radius guard of $10^{-10}$, and proposed
$\Rstar(n)\sim n^2/8$. We refer to that paper for the finite tables, search
algorithm and geometric regime analysis. The asymptotic theory here neither
uses those finite computations as a premise nor enlarges their certified
range.

Our main result is an effective characterization of the true global leading
coefficient. Sections~\ref{sec:line}--\ref{sec:lp} give its proof directly
from the angular model, through an all-pair line problem, genuine-label
concatenation and balanced finite-word programs. The lower and upper bounds
converge to the same constant with a prescribed error. Separate explicit
constructions give the interval stated in the abstract and disprove the
coefficient conjecture of \cite{v1}. Their complete proof dependencies are
specified in the supplement; the main variational characterization does
not assume either endpoint mechanism is optimal.

The line problem has a familiar geometric interpretation. Disks of radii
$a_i$ tangent to the horizontal axis, with centers $(2x_i,a_i)$, are
nonoverlapping precisely when $|x_i-x_j|\ge\sqrt{a_i a_j}$.
For $a_i\le1$, their horizontal envelope span differs from
$2(\max_i x_i-\min_i x_i)$ by at most two. Minimum-span disks on a shelf
are studied by Alt et al.\ \cite{shelf}. We do not claim this underlying
model as new; our asymptotic reduction, uniform-label recovery and effective
balanced-word characterization are proved directly below, without importing
an approximation or complexity theorem for that model.

\section{Exact angular comparison}\label{sec:line}

The cosine law gives the required pair angle
\begin{equation}\label{eq:theta}
 \theta_R(u,v)=2\asin\sqrt{\frac{uv}{(R+u)(R+v)}}
 =2\atan\sqrt{\frac{uv}{R(R+u+v)}}.
\end{equation}
Both directed cyclic arcs between a pair must be at least this angle.
This is equivalent to non-overlap, since their smaller separation belongs
to $[0,\pi]$. The angle is positive, symmetric, increasing in either
outer radius and strictly decreasing in $R>0$.
For a cyclic order $\sigma=(u_1,\ldots,u_n)$, $n\ge3$, with $u_{n+1}=u_1$,
$\ch(\sigma)$ is the unique positive root of
$\sum_{i=1}^n\theta_R(u_i,u_{i+1})=2\pi$.
Indeed, the sum decreases continuously from $n\pi$ to zero. Summing the
adjacent constraints in any full placement proves $\ch(\sigma)\le\full(\sigma)$.

\begin{definition}
For a finite multiset $a=(a_1,\ldots,a_N)$ of marks in $[0,1]$, let $b(a)$
be the least span of a line placement satisfying
\begin{equation}\label{eq:line}
 |x_i-x_j|\ge\sqrt{a_i a_j}\quad(i\ne j).
\end{equation}
Distinct zero-mark points may coincide. Write $b_n=b(1/n,2/n,\ldots,1)$.
\end{definition}

For a fixed word $w=(w_1,\ldots,w_N)$ in left-to-right order, its minimum
span $\ell(w)$ is attained by
\begin{equation}\label{eq:recurrence}
 x_1=0,\qquad x_i=\max_{j<i}\{x_j+\sqrt{w_iw_j}\},\qquad \ell(w)=x_N.
\end{equation}
Induction proves that these coordinates are componentwise no larger than
any feasible ordered coordinates with $x_1=0$; they also satisfy every
constraint. Nonnegative weights ensure ordered coordinates. Equivalently,
$\ell(w)$ is the longest increasing-index path from $1$ to $N$, with
edge weight $\sqrt{w_iw_j}$. Thus $b(a)$ is an attained minimum over a
finite set of orders. Deleting or decreasing marks cannot increase it;
scaling all marks by $h$ scales the minimum span by $h$.

\begin{lemma}[Bounded closing cost]\label{lem:closing}
Let $B(a)$ be the infimum circumference with both pair arcs at least
$\sqrt{a_i a_j}$. Then $b(a)\le B(a)\le b(a)+1$.
\end{lemma}
\begin{proof}
Cut a feasible circle to obtain a feasible line placement. Conversely,
close a feasible line placement by an additional gap of length one. One
arc for every pair is its line separation; the other contains the added
gap and is at least $1\ge\sqrt{a_i a_j}$. Both arcs are controlled.
\end{proof}

\begin{lemma}[Exact geometric squeeze]\label{lem:squeeze}
For every $n\ge3$,
\begin{equation}\label{eq:squeeze}
 \frac{b_n}{\pi n}-\frac1n
 \le\frac{\Rstar(n)}{n^2}\le\frac{b_n+1}{\pi n}.
\end{equation}
\end{lemma}
\begin{proof}
For $1\le u,v\le n$, the two expressions in \eqref{eq:theta} give
\[
 \frac{2\sqrt{uv}}{R+n}\le\theta_R(u,v)\le\frac{2\sqrt{uv}}R.
\]
The first uses $\asin z\ge z$ and
$\sqrt{(R+u)(R+v)}\le R+n$; the second uses $\atan z\le z$.
Scaling a feasible ring's angles by $(R+n)/(2n)$ gives a feasible linear
circle of circumference $\pi(R+n)/n$, proving the lower inequality.
Conversely, the circumference $b_n+1$ from Lemma~\ref{lem:closing},
scaled to $2\pi$, satisfies every original pair angle at
$R=n(b_n+1)/\pi$. This proves the upper inequality. No Taylor expansion,
adjacency reduction or assumption on the optimal order was used.
\end{proof}

\section{Existence by genuine-label concatenation}

\begin{theorem}[The leading limit exists]\label{thm:existence}
There is a constant $\CC$ such that
\begin{equation}\label{eq:limit}
 \Rstar(n)=\CC n^2+o(n^2),\qquad
 \pi\CC=E:=\lim_{n\to\infty}\frac{b_n}{n}
 =\inf_{n\ge2}\frac{b_n+1}{n}.
\end{equation}
\end{theorem}
\begin{proof}
Copies of a feasible line template, separated by unit gaps, satisfy all
cross-copy constraints as well as the internal ones. Fix $n\ge2$. Given
$N\ge n$, let $q=\lceil N/n\rceil$ and concatenate $q$ copies of an
optimal $b_n$ template. Its span is $q(b_n+1)-1$.
Retain the $N$ smallest marks. The $i$th sorted retained mark is
$\lceil i/q\rceil/n$. Scale all marks and coordinates by $h=nq/N$.
The resulting $i$th mark is at least $i/N$, so decreasing it to $i/N$
preserves every constraint. Each target label is assigned exactly once
to an actual retained point, with ties broken arbitrarily. Hence
\begin{equation}\label{eq:copies}
 b_N\le\frac{nq}{N}\bigl[q(b_n+1)-1\bigr].
\end{equation}
Intermediate marks exceeding one cause no problem: the line inequalities
are homogeneous, and final marks are exactly $1/N,\ldots,1$.

As $N\to\infty$, $q/N\to1/n$, giving
$\limsup_N b_N/N\le(b_n+1)/n$ for every fixed $n$.
The sequence is bounded between zero and one by equally spaced placements.
Taking $n$ along a subsequence attaining its liminf gives equality of
limsup and liminf. Taking the infimum and then $n\to\infty$ proves
the last equality in \eqref{eq:limit}. Lemma~\ref{lem:squeeze} transfers
the limit to the original ring problem for all sizes and both parities.
Finally, at least $n/2$ marks are at least $1/2$. Their successive line
separations are at least $1/2$, so $E\ge1/4$ and $\CC>0$.
\end{proof}

\section{Effective finite-word characterization}\label{sec:lp}

\subsection{Finite types and a uniform quantization error}
For $k\ge1$, let $A_{k,q}$ contain $q$ copies of each mark
$1/k,\ldots,1$, and let $A^-_{k,q}$ instead contain $q$ copies of
$0,1/k,\ldots,(k-1)/k$. Define
$e_k=\lim_{q\to\infty}b(A_{k,q})/(kq)$ and
$e^-_k=\lim_{q\to\infty}b(A^-_{k,q})/(kq)$; these limits exist.
Indeed $a_q=b(A_{k,q})+1$ is subadditive by unit-gap concatenation.
Writing $q=hp+s$ at fixed $p$ proves
$\limsup a_q/q\le a_p/p$; then taking $p$ along a liminf subsequence
proves the limit. The same argument applies to the lower types.

The two multisets share every nonzero type except the top type. Zero marks
cost nothing. Adding the $q$ unit marks as a separate equally spaced
block costs at most $q$, including the join. Thus, also for $k=1$,
\[
 e^-_k\le e_k\le e^-_k+\frac1k.
\]
For $n=kq$, the sorted uniform marks $i/n$ lie between the two sorted
type multisets. Theorem~\ref{thm:existence} and monotonicity imply
\begin{equation}\label{eq:types}
 e_k-\frac1k\le E\le e_k.
\end{equation}

\subsection{Balanced words and the quantified boundary relaxation}
Fix $k\ge1$ and $r\ge2$, and use all $k^r$ words
$w\in\{1/k,\ldots,1\}^r$, including repetitions. Let $c_j(w)$ count
occurrences of $j/k$. Compute $\ell(w)$ by \eqref{eq:recurrence}, and define
\begin{equation}\label{eq:lp}
\begin{split}
 \lambda_{k,r}=\min_p\ &\frac1r\sum_w p_w\ell(w),\\
 p_w\ge0,\quad &\sum_w p_w=1,\quad
 \sum_w p_wc_j(w)=\frac rk\quad(1\le j\le k).
\end{split}
\end{equation}
The polytope is nonempty: equally mix the $k$ constant words.

\begin{theorem}[Matching effective variational bounds]\label{thm:lp}
For every $k\ge1,r\ge2$,
\begin{equation}\label{eq:effective}
 \boxed{\quad\frac{\lambda_{k,r}-1/k}{\pi}\le\CC
 \le\frac{\lambda_{k,r}+1/r}{\pi}\quad}.
\end{equation}
Consequently $\CC=\pi^{-1}\lim_{k,r\to\infty}\lambda_{k,r}$, jointly,
with certified bracket width $(1/k+1/r)/\pi$.
\end{theorem}
\begin{proof}
Partition an optimal ordering of $A_{k,q}$ into consecutive blocks of
$r$ vertices, discarding fewer than $r$ leftover vertices. Their spans
sum to at most the full span and each is at least its word cost.
Empirical word distributions have convergent subsequences in the finite
simplex. Their limiting mean type counts are $r/k$, since the discarded
counts vanish after normalization. Hence $\lambda_{k,r}\le e_k$.

Conversely, the rational feasible polytope in \eqref{eq:lp} has a minimizing
rational vertex even though its cost entries are algebraic. Choose a
common denominator $D$. A batch with $Dp_w$ copies of each word has
exactly $Dr/k$ copies of each type, an integer by the count constraints.
Place every word at its minimum span and separate blocks by unit gaps.
All within-block and cross-block pairs are feasible. Repeating batches
gives asymptotic mean span
$\sum_w p_w(\ell(w)+1)/r=\lambda_{k,r}+1/r$.
The finite-type limit therefore satisfies
$e_k\le\lambda_{k,r}+1/r$. Combine with \eqref{eq:types}.
\end{proof}

\subsection{Exact certificates and computational meaning}
Directed rational square-root enclosures of width $\epsilon$ give
word-span enclosures of width at most $(r-1)\epsilon$: every path in
\eqref{eq:recurrence} has at most $r-1$ edges. Rational endpoint LPs
then enclose $\lambda_{k,r}$ with additional width at most
$(r-1)\epsilon/r$. They can be solved by finite exact vertex enumeration.
Practical solvers may instead discover candidates, followed by exact checks.
For a rational lower cost $\ell^-(w)$, dual coefficients satisfying
\[
 y_0+\sum_j y_jc_j(w)\le\ell^-(w)/r\quad\hbox{for every word}
\]
certify the lower value $y_0+(r/k)\sum_jy_j$. Feasible rational $p$
certifies an upper value using upper costs. All such checks are finite
rational comparisons; equality is not decided by indefinite interval
refinement. Together with a rational enclosure of $\pi$, this proves
arbitrary-precision computability of $\CC$ in principle.

The word count is exponential in $r$. No claim of cheap high-precision
evaluation or a closed-form constant follows. The bounded implementation
uses SciPy only to propose sparse primal and dual candidates; exact Fraction
arithmetic validates every equality and word inequality. Small word programs
need not improve the stronger explicit endpoint bounds below.

\section{Explicit geometric upper constructions}\label{sec:upper}

Fix $0\le\alpha<1$ and nonnegative adjacent slab lengths
$\ell_1,\ldots,\ell_k$ of total $T<1-\alpha$. Put $A=1+\alpha$,
$t_0=0$, $t_j=\sum_{h\le j}\ell_h$. On each slab
$[t_{j-1},t_j]$, use equal masses at
\[
 (t,A+t,A+t_{j-1}+t_j-t),\qquad
 (t,A+t_{j-1}+t_j-t,A+t).
\]
After $T$, use the diagonal $(t,1+\{t+\alpha\},1+\{t+\alpha\})$.
This defines a measure $\mu$ on $[0,1]\times[1,2]^2$ with uniform low
and separate high marginals. Its complete cell cost is
\begin{equation}\label{eq:blockcost}
 C(\mu)=\frac1{4\pi}\int
 \max\{\sqrt t(\sqrt X+\sqrt Y),\sqrt{XY}\}\,d\mu.
\end{equation}
Correct marginals alone do not establish finite recovery.

\begin{theorem}[General block transfer]\label{thm:blocks}
Every such finite partition has genuine alternating-half permutations with
all-pairs-feasible radii $\rho_m=\full(\sigma_m)$ and
\[
 \left|\frac{\rho_m}{(2m)^2}-C(\mu)\right|
 \le\frac{(6k^2+28k+1060)/m+16384/(3m^2)}{4\pi}
\]
for $m\ge\max\{2048,\lceil2/(1-\alpha-T)\rceil\}$.
It follows that $\CC\le C(\mu)$. Countably many slabs also transfer
if their total length is strictly below $1-\alpha$.
\end{theorem}
\begin{proof}[Proof with explicit construction]
Let $s=\lfloor\alpha m\rfloor$, $l_j=2\lfloor\ell_jm/2\rfloor$ and
$a_j=\sum_{h\le j}l_h$. Reverse only the even ranks within each
$(a_{j-1},a_j]$, fixing odd ranks, to obtain a permutation $J$.
Set $P_i=m+1+((J(i)+s-1)\bmod m)$ and
$\sigma_m=(1,P_1,\ldots,m,P_m)$, with $P_0=P_m$.
Disjoint involutions and the cyclic shift prove bijectivity, including
zero blocks and length-two identities. The integer gate keeps the final
block endpoint at least two ranks before the high wrap, and its exit
strictly before that wrap.

Double panels of width $2/m$, allocated half to each parity, recover
the two reflected orientations. Interior coordinate error is at most
$4/m$. There are at most $k+3$ exceptional actual cells: shared exits,
the cyclic predecessor, and the two wrap cells. Rounded slab boundaries
and wrap differ on a set of length at most $[k(k+1)+1]/m$, with remaining
coordinate error at most $(4k-1)/m$. Write
$g(t,X,Y)=\max\{\sqrt t(\sqrt X+\sqrt Y),\sqrt{XY}\}$ for the integrand
in \eqref{eq:blockcost}. It is $4$-Lipschitz and lies in $[1,3)$.
The actual normalized full-cell moment is
\[
 G_m=\frac1m\sum_{i=1}^m g(i/m,P_{i-1}/m,P_i/m).
\]
It differs from
$\int g\,d\mu$ by at most $(6k^2+28k+36)/m$.

For arbitrary distinct highs in $[m+1,2m]$, the exact full-feasibility
criterion for this alternating order is
\[
 \sum_{i=1}^m\max\{\theta_R(P_{i-1},i)+\theta_R(i,P_i),
                         \theta_R(P_{i-1},P_i)\}\le2\pi.
\]
Its sufficiency follows from the high-shell triangle inequality
$2\theta_R(m+1,m+1)>\theta_R(2m,2m)$ and high/low separation:
each directed high path contracts through whole cells, and paths with
low endpoints either contain such a high path or an incident high gap
that already exceeds the endpoint requirement. Both directions are
treated separately. At its unique root, choose the first half of each
cell at its required low/high angle and put any excess in the other half.
This closes at $2\pi$ and satisfies all pairs.

At $R=4cm^2$, $c\ge1/32$, the full score differs from $G_m/(2c)$ by
at most $1024/m+16384/(3m^2)$. For $m\ge2048$ this first brackets
the actual normalized root in $(1/32,1/2)$, then yields the stated error.
Deleting label $2m$ gives the odd upper bound. For countably many slabs,
truncate and use that the omitted complete-cost integral is at most six
times the omitted length. A diagonal sequence with $m$ growing faster
than $k^3$ pays both finite and tail errors. Full cell/seam formulas and
the angle estimate are supplied in the versioned proof supplement \cite{blocks}.
\end{proof}

For the retained explicit construction, let $x_*$ be the unique minimizer
on $[0,1]$ of
\[
 E(x)=\int_0^x\!\left[
 \max\{\sqrt{(1+t)(1+x-t)},\sqrt t(\sqrt{1+t}+\sqrt{1+x-t})\}
 -\max\{1+t,2\sqrt{t(1+t)}\}\right]dt.
\]
Let $K(\alpha)=\int_0^1\max\{\sqrt{t h_\alpha(t)},h_\alpha(t)/2\}\,dt$,
where $h_\alpha(t)=1+\{t+\alpha\}$, and let $\widehat\alpha$ be the
unique zero of $K'(\alpha)+(1+\alpha)E(x_*)$ on $[0,1/2]$.
Set $A=1+\widehat\alpha$ and $\lambda=Ax_*$.
The second width $\epsilon_b$ is the unique mixed-branch stationary
minimum for a separate reflection starting at $\lambda$; the third width
$\Delta_*$ is the analogous unique minimum starting at
$v=\lambda+\epsilon_b$. Each width is minimized over
$[0,A/3-u]$ at its fixed start $u$, retaining the full max.
The supplement \cite{upperinputs} gives their exact defining integrals,
existence/uniqueness proofs and rational isolations:
\[
\begin{gathered}
 .1093<\widehat\alpha<.10931,\qquad .2876<x_*<.2877,\\
 .043<\epsilon_b<7/160,\qquad 29/5000<\Delta_*<27/4000.
\end{gathered}
\]
These terminating decimals are rational bracket endpoints, not definitions
or numerical replacements for the parameters.

Let $C_3(d)$ be \eqref{eq:blockcost} for the three lengths
$(\lambda,\epsilon_b,d)$, and $C_4(\eta)$ the same cost for the four
lengths $(\lambda,\epsilon_b,\Delta_*,\eta)$.
Use $\eta=1/20000$ and define $U_4=C_4(1/20000)$.
If $w=\lambda+\epsilon_b+\Delta_*$, $B=A+w$ and
$M=B+\eta/2$, the fourth increment at $0<\eta\le1/20000$ is exactly
\[
 C_4(\eta)-C_3(\Delta_*)=
 -\frac1{4\pi}\int_{-\eta/2}^{\eta/2}
       \frac{z^2}{M+\sqrt{M^2-z^2}}\,dz.
\]
Rational input margins keep both changed maxima strictly chord.
Since $2B\le M+\sqrt{M^2-z^2}\le2M$, $M<3/2$ and $\pi<4$,
\begin{equation}\label{eq:U4}
 \CC\le U_4<C_3(\Delta_*)-\frac1{4608000000000000}.
\end{equation}
The leading decrement is $-\eta^3/(96\pi B)$, with remainder between
zero and $\eta^4/(192\pi B^2)$. The theorem transfers this literal
continuous saving to genuine full geometry. It does not make four blocks
globally optimal or optimize an infinite partition.

\section{An explicit global lower endpoint}

Define $\tau$ as the unique root of $\cos\tau=\tau$ in $(0,1)$, set
$q=(1-\sin\tau)/(1+\sin\tau)$ and
$C_{\rm term}=\tau(1+q)/(2\pi)$.
For the fixed cutoffs
\[
 (\beta_1,\beta_2,\beta_3)=(2091/10000,10907/50000,23/100)
\]
put
\[
 D_i=\int_q^{\beta_i}[1+q-x-2\sqrt{x(1+q-x)}]\,dx,
 \quad F(h)=\sum_i h_iD_i-8\sum_i h_i^3.
\]
Let $a=113/12500$, $b=593/50000$. On
$h_i\ge0$, $h_1+h_2\le a$, $h_2+h_3\le b$, define
\[
 \eta_{\rm width}=\max_h
 \frac{\max\{F(h),0\}}{\pi(16+432\sum_i h_i)}.
\]
This maximum is uniquely attained at $(a-z_*,z_*,b-z_*)$, where $z_*$
is the unique zero on $[0,a]$ of $G=N'L+432N$,
$N(z)=F(a-z,z,b-z)$ and $L(z)=16+432(a+b-z)$.

\begin{proposition}[Explicit global lower bound]\label{prop:lower}
The global constant satisfies $\CC\ge C_{\rm term}+\eta_{\rm width}$,
with exact directed enclosure
\[
 .14056946887766098063257
 < C_{\rm term}+\eta_{\rm width}
 < .14056946887766098063392.
\]
In particular $\CC>1/8$, so both the v1 coefficient conjecture and its
stronger $n^2/8-\Rstar(n)=O(\sqrt n)$ claim are false.
\end{proposition}
\begin{proof}[Dependency-explicit proof outline]
The complete proof and exact arithmetic certificates are in \cite{lower}.
Its essential coupled information is one shared crossing energy for the
same order and all three nested terminal restrictions. An explicit assignment
dual bounds the reflected-grid energy by eight times the excess chain cost.
The absolute signed deletion discrepancy at cutoff $i$ is at most $54$
times this energy plus twice its crossing term. Separated strips charge
the common energy once, yielding
the scalar gain $\max\{F(h),0\}/(16+432\sum h_i)$.
The exact angular squeeze transfers each induced chain bound to any original
full feasible configuration. Uniform finite errors vanish at every fixed
strictly separated width vector.

For the width optimum, $G'<0$ and rational endpoint gates isolate its zero.
Set $u=(a-z_*,z_*,b-z_*)$, $\mathcal L(h)=16+432\sum_i h_i$ and
$\rho=F(u)/\mathcal L(u)$. The gradient of $F-\rho\mathcal L$ at $u$ is a positive
combination of the two constraint normals. Its exact remainder is
$-8\sum_i(h_i-u_i)^2(h_i+2u_i)$, proving global uniqueness on the whole
closed region. To pass to the geometric limit, first use $tu$ at fixed
$0<t<1$, take $n\to\infty$, then $t\uparrow1$. The boundary
widths fail finite floor separation infinitely often; the proof never uses
them as an eventual finite gate. The uniform finite error bounds for this
limiting argument are supplied in \cite{lower}. The result does not assert
that this particular lower mechanism reaches $\CC$.
\end{proof}

\section{Reproducibility and open questions}

The supplement \cite{supplement} is fixed at commit
\texttt{3beb8d70c5b3748d370a92855847bdf574e5a14f}.
All supplementary references below use that snapshot. It contains detailed
proofs of the explicit endpoint estimates and block-transfer bounds, together
with standalone checkers for line placements, label recovery, reflected
cells and rational primal/dual word certificates. Directed arithmetic checks
include deliberately omitted nonadjacent constraints, missing closing gaps,
invalid label assignments and corrupt LP certificates. These bounded checks
corroborate finite computations; the analytic arguments establish the infinite
quantifiers in the theorems.

The leading normalized behavior is characterized effectively, while neither
explicit endpoint is known to equal $\CC$. Practical evaluation or
simplification of the constant, sharp explicit constructions, microscopic
optimizer structure and subleading terms remain open. The number of words
in the finite programs is exponential in their length; computability does
not give an efficient high-precision algorithm. The limit theorem neither
certifies a finite optimum beyond $n=14$ nor establishes a universal
floating-circle cascade or a contact structure for global optimizers.

\section*{Acknowledgment}
AI assistance was used for proof exploration, software, drafting and
adversarial checks. Responsibility for the mathematical statements and
presentation rests with the author.

\begin{thebibliography}{9}
\small\raggedright
\bibitem{v1} M. Falconi, \emph{Arranging circles of radii $1,2,\ldots,n$
around a central circle: a Supnick TSP and certified finite optima},
arXiv:2607.28654, version 1 (2026).
\url{https://arxiv.org/abs/2607.28654v1}.
\bibitem{shelf} H. Alt, K. Buchin, S. Chaplick, O. Cheong, P. Kindermann,
C. Knauer and F. Stehn, \emph{Placing your coins on a shelf}, Journal of
Computational Geometry 9(1), 312--327 (2018).
\url{https://doi.org/10.20382/jocg.v9i1a10}.
\bibitem{blocks} M. Falconi, Ringmin proof supplement,
\nolinkurl{research/PERMUTED_ALTERNATING_HALVES.md}, Sections 1--6;
\nolinkurl{research/PERMUTED_HALVES_GENERAL_BLOCK_TRANSFER.md};
\nolinkurl{research/PERMUTED_HALVES_FOURTH_ADJACENT_BLOCK.md}.
Snapshot specified in \cite{supplement}.
\bibitem{upperinputs} M. Falconi, Ringmin proof supplement,
\nolinkurl{research/PERMUTED_HALVES_REFLECTED_PREFIX_LAMBDA.md}, Sections 2--7;
\nolinkurl{research/PERMUTED_HALVES_REFLECTED_PREFIX_ALPHA_MINIMUM.md}, Sections 1, 3--4;
\nolinkurl{research/PERMUTED_HALVES_SECOND_BLOCK_BOUNDARY_MINIMUM.md};
\nolinkurl{research/PERMUTED_HALVES_THIRD_BLOCK_MIXED_WIDTH.md}.
Snapshot specified in \cite{supplement}.
\bibitem{lower} M. Falconi, Ringmin proof supplement,
\nolinkurl{research/COMMON_CHAIN_QUANTITATIVE_STABILITY.md}, Sections 2--5, 11;
\nolinkurl{research/THREE_LEVEL_COMMON_CHAIN.md}, Sections 7, 9, 11--12.
Snapshot specified in \cite{supplement}.
\bibitem{supplement} M. Falconi, Ringmin mathematical and computational
supplement, commit \texttt{3beb8d70c5b3748d370a92855847bdf574e5a14f}.
\href{https://github.com/falker47/ringmin/tree/3beb8d70c5b3748d370a92855847bdf574e5a14f}{Repository snapshot}.
Main variational proof:
\nolinkurl{research/GLOBAL_ASYMPTOTIC_VARIATIONAL_LIMIT.md}.
\end{thebibliography}
\end{document}
